\scanpage{71}
\uline{Th} Si un syst. de $\rho_{\vec D,s}$ existe satisfaisant
a) b) c) d) e) f) \uline{et si} il n'existe pas de ps. droites
% slanted
\margin{il suffit de \deleted{d)} e), f) \unsure{etc|et} (ss doute) a) en est une \unsure{conséqu.|conséquence}}
qui sont des monogônes ou bigône, alors
ce système des $\rho_{\vec D,s}$ est unique.

De plus, le syst. des $\Omega(s)=\{\rho_s,\rho_s^{-1}\}$ -- est
\uline{de façon unique} un syst. local satisfaisant
\deleted{\unsure{si|Si} \illegible{1}} \ldots.
% slanted
\margin{Pour pouvoir donner un sens au syst. local, il faut de plus \uline{\unsure{poser|fixer|fortes}} comme hyp. suppl. \uline{b)} \deleted{\unsure{et d)}} c) \deleted{pour pouvoir \unsure{les}}}
% archive sheet number: 70
